\documentclass[10pt,a4paper]{amsart}
\usepackage[margin=19mm]{geometry}
\usepackage{amsmath,amssymb,mathtools}
\usepackage[scr=rsfs,cal=euler]{mathalfa}
\usepackage{xfrac}
\usepackage[unicode,hidelinks,bookmarks=false]{hyperref}
\newcommand{\ie}{i.e.\ }
\newcommand{\cf}{cf.\ }
\newcommand{\resp}{respectively\ }
\newcommand{\Subsection}{\subsection}
\providecommand{\nobreakdash}{\nobreak\hbox{-}}
\setlength{\emergencystretch}{2em}
\multlinegap0pt
\usepackage{amssymb}
\usepackage{xcolor}
\usepackage{float}
\usepackage{tikz}
\usepackage{mathtools}
\newtheorem{theorem}{Theorem}
\renewcommand{\Im}{\operatorname{Im}}%imaginary part
\renewcommand{\Re}{\operatorname{Re}}%real part
\title{Eigenvalues of the normalized complex Laplacian on finite electrical networks}
\author{Anna Muranova*, Robert Schippa}
\date{Source: arXiv:2012.12759v1 (CC BY 4.0)}
\begin{document}
\maketitle

\noindent\emph{Source and licence.} Eigenvalues of the normalized complex Laplacian on finite electrical networks. Version arXiv:2012.12759v1 (CC BY 4.0), distributed by arXiv under the Creative Commons Attribution 4.0 International licence, http://creativecommons.org/licenses/by/4.0/, as recorded in the arXiv licence metadata for this exact version and shown on the abstract page https://arxiv.org/abs/2012.12759v1. Full licence notice: \texttt{licenses/arxiv-2012.12759-CC-BY-4.0.txt}.\par\smallskip

\noindent\textit{Editorial scope.} Counted unit: the article's theorem thm:CircleRegions (source0.tex lines 144--155). The article's complete original proof of it is delivered with it (source0.tex lines 381--486, one \texttt{proof} environment of 106 lines, which the article places away from the statement). No mathematics is written, repaired or re-derived: the statement and the proof are the article's own lines, in the article's own notation, and no retained line is substituted or re-flowed.  Retained uncounted as the source-local context the proof needs: lines 210--213.  Nothing was omitted from any retained span, so the package carries no omission record.  No retained line contains a citation, so the wrapper carries no bibliography and no bibliography entry.  No retained line contains a figure environment, an image-inclusion command or drawing code, so no image is delivered and the package declares no asset dependency.  Editorial formatting only, in addition: an AMS \texttt{amsart} preamble that restates the article's own declarations and macros used by the retained lines, namely the environments \texttt{theorem} on separate counters, together with the standard packages those lines need.  The retained environments print in this document as Theorem 1 in order of appearance, whereas the article prints the same items with the numbering of its own full document.  The complete native source of the article as distributed in the arXiv e-print for this exact version is bundled unchanged as \texttt{sources/106216/source0.tex}.
\begin{equation}
\label{eq:EstimateModulusAdmittance}
| \rho_{xy} | \leq \frac{|s|}{\Re s} \Re \rho_{xy}.
\end{equation}

\begin{theorem}
\label{thm:CircleRegions}
Let $\Delta_\rho$ be a normalized complex-weighted Laplacian. Then the following holds:
\begin{itemize}
\item
All its eigenvalues with positive real part lie in the circle with center at $(1,|\Im s|/\Re s)$ and radius $\sqrt {1+(\Im s/\Re s)^2}$.
\item
All its eigenvalues with negative real part lie in the circle with the center at $(1,-|\Im s|/\Re s)$ and radius $\sqrt {1+(\Im s/\Re s)^2}$.
\item
All its real eigenvalues lie in $[0,2]$.
\end{itemize}
\end{theorem}

\begin{proof}[Proof of Theorem \ref{thm:CircleRegions}]
Let ${\lambda=u+iw}$ be an eigenvalue of the normalized complex-weighted Laplacian $\widetilde \Delta_\rho$ with the eigenfunction $f$, i.e, 
\begin{equation*}
\lambda f(x)=\widetilde \Delta_\rho f(x).
\end{equation*}
Then by Green's formula we have
\begin{equation}\label{eveq}
\sum_{x\in V}\lambda|f(x)|^2\rho(x)=\dfrac{1}{2}\sum_{x,y\in V}|\nabla_{xy}f|^2\rho_{xy}.
\end{equation}
Let us write separately real and imaginary parts of the last equality, assuming  ${\rho_{xy}=\tau_{xy}+i\sigma_{xy}}$ and ${\rho(x)=\tau(x)+i\sigma(x)}$. Note that in this case $\tau(x)=\sum_{y}\tau_{xy}$ and $\sigma(x)=\sum_{y}\sigma_{xy}$. 

Then,
\begin{equation*}
\lambda\sum_{x\in V}|f(x)|^2\rho(x)=(u+iw)\sum_{x\in V}|f(x)|^2(\tau(x)+i\sigma(x)),
\end{equation*}
and
\begin{equation*}
\dfrac{1}{2}\sum_{x,y\in V}|\nabla_{xy}f|^2\rho_{xy}=\dfrac{1}{2}\sum_{x,y\in V}|\nabla_{xy}f|^2(\tau_{xy}+i\sigma_{xy}).
\end{equation*}


Therefore, for the real part of \eqref{eveq}, we have
\begin{equation}\label{realpart}
u\sum_{x\in V}|f(x)|^2\tau(x)-w\sum_{x\in V}|f(x)|^2\sigma(x)=\dfrac{1}{2}\sum_{x,y\in V}|\nabla_{xy}f|^2\tau_{xy},
\end{equation}
and for the imaginary part,
\begin{equation}\label{imaginarypart}
u\sum_{x\in V}|f(x)|^2\sigma(x)+w\sum_{x\in V}|f(x)|^2\tau(x)=\dfrac{1}{2}\sum_{x,y\in V}|\nabla_{xy}f|^2\sigma_{xy}.
\end{equation}
Multiplying \eqref{realpart} by $u$ and \eqref{imaginarypart} by $w$ and summing the obtained equalities up, we get

\begin{equation}\label{eq1}
(u^2+w^2)\sum_{x\in V}|f(x)|^2\tau(x)=\dfrac{1}{2}\sum_{x,y\in V}|\nabla_{xy}f|^2(\tau_{xy}u+\sigma_{xy}w).
\end{equation}
Let us estimate the left-hand side, using the obvious inequality 
\begin{equation*}
|\nabla_{xy}f|^2=|f(y)-f(x)|^2\le 2(|f(y)|^2+|f(x)|^2).
\end{equation*}
%which is true, due the following computations
%\begin{align*}
%|a-b|^2&=(a-b)(\overline a-\overline b)=|a|^2-a\overline b-\overline a b+|b|^2\\
%=&-|a|^2-a\overline b-\overline a b-|b|^2+2(|a|^2+|b|^2)\\
%=&-|a+b|^2+2(|a|^2+|b|^2)\le2(|a|^2+|b|^2).
%\end{align*}
Moreover, the right-hand side of \eqref{eq1} is positive, this means also the left-hand side is positive and at least one $(\tau_{xy}u+\sigma_{xy}w)$ is positive.

Let $w>0$.
Then, we have
\begin{align*}
(u^2+w^2)\sum_{x\in V}|f(x)|^2\tau(x)&=\dfrac{1}{2}\sum_{x,y\in V}|\nabla_{xy}f|^2\tau_{xy}(u+\dfrac{\sigma_{xy}}{\tau_{xy}}w)\\
&\le \sum_{x,y\in V}(|f(y)|^2+|f(x)|^2)\tau_{xy}(u+\dfrac{\sigma_{xy}}{\tau_{xy}}w) \\
&\le \sum_{x,y\in V}(|f(y)|^2+|f(x)|^2)\tau_{xy}(u+w\max_{x\sim y}\dfrac{\sigma_{xy}}{\tau_{xy}}) \\
&\le 2\sum_{x\in V}|f(x)|^2\tau(x)(u+w\max_{x\sim y}\dfrac{\sigma_{xy}}{\tau_{xy}}) .
\end{align*}
Therefore,
\begin{align*}
u^2+w^2&\le 2(u+w\max_{x\sim y}\dfrac{\sigma_{xy}}{\tau_{xy}}) ,
\end{align*}
i.e.,
\begin{align*}
(u-1)^2+\left(w-\max_{x\sim y}\dfrac{\sigma_{xy}}{\tau_{xy}}\right)^2&\le 1+\left(\max_{x\sim y}\dfrac{\sigma_{xy}}{\tau_{xy}}\right)^2.
\end{align*}

Let $w<0$.
Then, we have
\begin{align*}
(u^2+w^2)\sum_{x\in V}|f(x)|^2\tau(x)&=\dfrac{1}{2}\sum_{x,y\in V}|\nabla_{xy}f|^2\tau_{xy}(u+\dfrac{\sigma_{xy}}{\tau_{xy}}w)\\
&\le \sum_{x,y\in V}(|f(y)|^2+|f(x)|^2)\tau_{xy}(u+\dfrac{\sigma_{xy}}{\tau_{xy}}w) \\
&\le \sum_{x,y\in V}(|f(y)|^2+|f(x)|^2)\tau_{xy}(u+w\min_{x\sim y}\dfrac{\sigma_{xy}}{\tau_{xy}}) \\
&\le 2\sum_{x\in V}|f(x)|^2\tau(x)(u+w\min_{x\sim y}\dfrac{\sigma_{xy}}{\tau_{xy}}) .
\end{align*}
Therefore,
\begin{align*}
u^2+w^2&\le 2(u+w\min_{x\sim y}\dfrac{\sigma_{xy}}{\tau_{xy}}) ,
\end{align*}
i.e.,
\begin{align*}
(u-1)^2+\left(w-\min_{x\sim y}\dfrac{\sigma_{xy}}{\tau_{xy}}\right)^2&\le 1+\left(\min_{x\sim y}\dfrac{\sigma_{xy}}{\tau_{xy}}\right)^2.
\end{align*}

Let $w=0$.
Then, we have
\begin{align*}
u^2\sum_{x\in V}|f(x)|^2\tau(x)&=\dfrac{1}{2}\sum_{x,y\in V}|\nabla_{xy}f|^2\tau_{xy}u\\
&\le 2\sum_{x\in V}|f(x)|^2\tau(x)u .
\end{align*}
Therefore,
\begin{align*}
u^2&\le 2u,
\end{align*}
which means that the real eigenvalues of the normalized complex Laplacian lie in $[0,2]$.

From \eqref{eq:EstimateModulusAdmittance} follows by squaring
\begin{equation*}
\dfrac{\Re^2 \rho_{xy}}{\Re^2 \rho_{xy}+\Im^2 \rho_{xy}}\ge \dfrac{\Re^2 s}{\Re^2 s+\Im^2 s}.
\end{equation*}
This yields
\begin{equation*}
\dfrac{\Im^2 \rho_{xy}}{\Re^2 \rho_{xy}}\le \dfrac{\Im^2 s}{\Re^2 s},
\end{equation*}
and hence,
\begin{equation*}
-\left| \dfrac{\Im s}{\Re s}\right|\le \dfrac{\Im \rho_{xy}}{\Re\rho_{xy}}\le\left| \dfrac{\Im s}{\Re s}\right|.
\end{equation*}
The proof is complete.
\end{proof}

\end{document}
